% Balanced complete-edition layout; entries unchanged from the accepted bibliography.
\begin{thebibliography}{99}
\label{refs:chapter11}

\bibitem[Aral et~al.(2009)]{aral2009}Aral, S., Muchnik, L., and Sundararajan, A. (2009).
Distinguishing influence-based contagion from homophily-driven diffusion in dynamic networks.
\emph{Proceedings of the National Academy of Sciences}, 106(51), 21544--21549.
\doi{10.1073/pnas.0908800106}.

\bibitem[Benkler(2002)]{benkler2002}Benkler, Y. (2002).
Coase's penguin, or, Linux and the nature of the firm.
\emph{Yale Law Journal}, 112, 369--446.
\url{https://yalelawjournal.org/article/coases-penguin-or-linux-and-the-nature-of-the-firm}.

\bibitem[Borgatti(2005)]{borgatti2005}Borgatti, S.~P. (2005).
Centrality and network flow.
\emph{Social Networks}, 27(1), 55--71.
\doi{10.1016/j.socnet.2004.11.008}.

\bibitem[Borsboom et~al.(2004)]{borsboom2004}Borsboom, D., Mellenbergh, G.~J., and van Heerden, J. (2004).
The concept of validity.
\emph{Psychological Review}, 111(4), 1061--1071.
\doi{10.1037/0033-295X.111.4.1061}.

\bibitem[Due\~nas et~al.(2021)]{duenas2021}Due\~nas, S., Cosentino, V., Gonzalez-Barahona, J.~M., et~al. (2021).
GrimoireLab: A toolset for software development analytics.
\emph{PeerJ Computer Science}, 7, e601.
\doi{10.7717/peerj-cs.601}.

\bibitem[Espeland and Sauder(2007)]{espeland2007}Espeland, W.~N., and Sauder, M. (2007).
Rankings and reactivity: How public measures recreate social worlds.
\emph{American Journal of Sociology}, 113(1), 1--40.
\doi{10.1086/517897}.

\bibitem[Geiger and Ribes(2011)]{geiger2011}Geiger, R.~S., and Ribes, D. (2011).
Trace ethnography: Following coordination through documentary practices.
In \emph{Proceedings of the 44th Hawaii International Conference on System Sciences}, pp. 1--10.
\url{https://stuartgeiger.com/trace-ethnography-hicss-geiger-ribes.pdf}.

\bibitem[Han et~al.(2024)]{han2024ose}Han, F., Zhao, S., Wang, W., Zhou, A., Qian, W., Zhou, X., Peng, J., You, L., Chen, Y., Xia, X., Tang, Y., Yang, L., and Tian, C. (2024).
Open Source Evaluatology: An evaluation framework and methodology for open source ecosystems based on evaluatology.
\emph{BenchCouncil Transactions on Benchmarks, Standards and Evaluations}, 4(4), 100190.
\doi{10.1016/j.tbench.2025.100190}.
\emph{Assigned to the December 2024 issue; published online in 2025.}

\bibitem[Han et~al.(2025)]{han2025openperf}Han, F., Bi, F., Zhang, Y., Li, H., You, L., Wang, W., Qian, W., and Zhou, A. (2025).
OpenPerf: A data science benchmark system for open source ecosystem sustainable development.
\emph{Chinese Journal of Computers}, 48(3), 632--649. In Chinese.
\doi{10.11897/SP.J.1016.2025.00632}.

\bibitem[Hicks et~al.(2015)]{hicks2015}Hicks, D., Wouters, P., Waltman, L., de Rijcke, S., and Rafols, I. (2015).
Bibliometrics: The Leiden Manifesto for research metrics.
\emph{Nature}, 520, 429--431.
\doi{10.1038/520429a}.

\bibitem[Huang et~al.(2024)]{huang2024}Huang, W., Xia, X., Zhou, A., Zhou, X., Wang, W., Zhao, S., Wang, Z., and Bian, S. (2024).
OSGraph: A data visualization insight platform for open source community.
In \emph{Database Systems for Advanced Applications (DASFAA 2024)}, LNCS 14856, pp. 476--479. Springer.
\url{https://doi.org/10.1007/978-981-97-5575-2_42}.

\bibitem[Kitchenham et~al.(1995)]{kitchenham1995}Kitchenham, B., Pfleeger, S.~L., and Fenton, N. (1995).
Towards a framework for software measurement validation.
\emph{IEEE Transactions on Software Engineering}, 21(12), 929--943.
\doi{10.1109/32.489070}.

\bibitem[Lave and Wenger(1991)]{lave1991}Lave, J., and Wenger, E. (1991).
\emph{Situated Learning: Legitimate Peripheral Participation}. Cambridge University Press.
\doi{10.1017/CBO9780511815355}.

\bibitem[Messick(1995)]{messick1995}Messick, S. (1995).
Validity of psychological assessment: Validation of inferences from persons' responses and performances as scientific inquiry into score meaning.
\emph{American Psychologist}, 50(9), 741--749.
\doi{10.1037/0003-066X.50.9.741}.

\end{thebibliography}
\clearpage
\begingroup
\renewcommand{\bibsection}{\section*{References\enspace{\normalfont\sffamily\fontsize{11}{14}\selectfont continued}}}
\begin{thebibliography}{99}
\bibitem[Munzner(2009)]{munzner2009}Munzner, T. (2009).
A nested model for visualization design and validation.
\emph{IEEE Transactions on Visualization and Computer Graphics}, 15(6), 921--928.
\url{https://www.cs.ubc.ca/labs/imager/tr/2009/NestedModel/}.

\bibitem[Nicol and Macfarlane-Dick(2006)]{nicol2006}Nicol, D.~J., and Macfarlane-Dick, D. (2006).
Formative assessment and self-regulated learning: A model and seven principles of good feedback practice.
\emph{Studies in Higher Education}, 31(2), 199--218.
\doi{10.1080/03075070600572090}.

\bibitem[Ostrom(2009)]{ostrom2009}Ostrom, E. (2009).
A general framework for analyzing sustainability of social-ecological systems.
\emph{Science}, 325(5939), 419--422.
\doi{10.1126/science.1172133}.

\bibitem[Peng(2011)]{peng2011}Peng, R.~D. (2011).
Reproducible research in computational science.
\emph{Science}, 334(6060), 1226--1227.
\doi{10.1126/science.1213847}.

\bibitem[Star and Griesemer(1989)]{star1989}Star, S.~L., and Griesemer, J.~R. (1989).
Institutional ecology, `translations' and boundary objects: Amateurs and professionals in Berkeley's Museum of Vertebrate Zoology, 1907--39.
\emph{Social Studies of Science}, 19(3), 387--420.
\doi{10.1177/030631289019003001}.

\bibitem[Tang et~al.(2024)]{tang2024}Tang, Y., Zhao, S., Xia, X., Bi, F., and Wang, W. (2024).
HyperCRX: A browser extension for insights into GitHub projects and developers.
In \emph{Proceedings of the 32nd IEEE/ACM International Conference on Program Comprehension}, pp. 460--464. ACM.
\doi{10.1145/3643916.3644440}.

\bibitem[Wang et~al.(2025)]{wang2025hyper}Wang, Y., Wang, W., Xia, X., Tang, Y., Bi, F., Han, F., Peng, J., Zhao, S., and You, L. (2025).
HyperCRX 2.0: A comprehensive and automated tool for empowering GitHub insights.
In \emph{2025 IEEE/ACM 47th International Conference on Software Engineering: Companion Proceedings (ICSE-Companion)}, pp. 49--52.
\emph{Figures and implementation details follow the author-supplied paper.}

\bibitem[Wang et~al.(2026)]{wang2026digger}Wang, W., Han, F., Zhao, S., Zhou, X., Qian, W., Zhou, A., Xia, X., Yang, L., Wang, R., Jiang, N., and Duan, M. (2026).
OpenDigger: A practical framework for assessing community health and sustainability in open source collaboration platforms.
In \emph{Proceedings of the ACM Web Conference 2026}, pp. 7745--7756. ACM.
\doi{10.1145/3774904.3792807}.

\bibitem[Wilkinson et~al.(2016)]{wilkinson2016}Wilkinson, M.~D., Dumontier, M., Aalbersberg, I.~J., et~al. (2016).
The FAIR Guiding Principles for scientific data management and stewardship.
\emph{Scientific Data}, 3, 160018.
\doi{10.1038/sdata.2016.18}.

\bibitem[Wise and Shaffer(2015)]{wise2015}Wise, A.~F., and Shaffer, D.~W. (2015).
Why theory matters more than ever in the age of big data.
\emph{Journal of Learning Analytics}, 2(2), 5--13.
\doi{10.18608/jla.2015.22.2}.

\bibitem[X-lab(2024)]{xlab2024talent}X-lab (2024).
\emph{OpenTalent: An overview of open source in Chinese universities}.
Dashboard screenshot supplied by the authors; displayed update date December 2024. Data source identified in the interface: OpenDigger.

\bibitem[Zhan et~al.(2026)]{zhan2026book}Zhan, J., Wang, L., Gao, W., Wang, C., Fan, F., Kang, G., and Li, H. (2026).
\emph{Evaluatology: The Science of Revealing Causes and Effects} [title translated from Chinese].
Chinese edition, version 1.0, August 2026. BenchCouncil Press.
ISBN 978-988-71596-9-8. \doi{10.66834/33wi1xx8}.

\bibitem[Zhang et~al.(2024)]{zhang2024galaxy}Zhang, X., Zhao, S., Tang, Y., Xia, X., and Wang, W. (2024).
OpenGalaxy: An interactive exploration platform for a visualized GitHub full domain collaboration network.
In \emph{Proceedings of the 32nd IEEE/ACM International Conference on Program Comprehension}, pp. 455--459. ACM.
\doi{10.1145/3643916.3644441}.

\end{thebibliography}
\endgroup
